\documentclass[12pt,a4paper]{report}
\usepackage[utf8]{inputenc}
\usepackage[T1]{fontenc}
\usepackage[french]{babel}
\usepackage{geometry}
\geometry{a4paper, margin=2.4cm, top=2.6cm, bottom=2.6cm}
\usepackage{setspace}
\setstretch{1.38}
\usepackage{hyperref}
\usepackage{listings}
\usepackage{xcolor}
\usepackage{booktabs}
\usepackage{array}
\usepackage{amsmath}
\usepackage{mdframed}
\usepackage{fancyhdr}
\usepackage{titlesec}

\pagestyle{fancy}
\fancyhf{}
\rhead{\textbf{Rapport Technique : Power Query}}
\lhead{Synchronisation Excel Multi-Feuilles}
\rfoot{Page \thepage}

\definecolor{codegreen}{rgb}{0,0.6,0}
\definecolor{codegray}{rgb}{0.5,0.5,0.5}
\definecolor{codepurple}{rgb}{0.58,0,0.82}
\definecolor{backcolour}{rgb}{0.96,0.96,0.94}
\definecolor{primaryblue}{rgb}{0.1,0.3,0.6}

\lstdefinestyle{mystyle}{
    backgroundcolor=\color{backcolour},   
    commentstyle=\color{codegreen},
    keywordstyle=\color{primaryblue}\bfseries,
    numberstyle=\tiny\color{codegray},
    stringstyle=\color{codepurple},
    basicstyle=\ttfamily\footnotesize,
    breakatwhitespace=false,         
    breaklines=true,                 
    captionpos=b,                    
    keepspaces=true,                 
    numbers=left,                    
    numbersep=7pt,                  
    showspaces=false,                
    showstringspaces=false,
    showtabs=false,                  
    tabsize=2
}
\lstset{style=mystyle}

\title{\Huge \textbf{Manuel Technique et Guide d'Architecture d'Entreprise : Synchronisation Excel avec Power Query (Connexion Externe Locale)}\\[1.5em]
\Large Solution d'Ingénierie de Données Intra-Classeur pour la Duplication Dynamique de Feuille Maître vers Feuilles Copies}
\author{\textbf{Département des Systèmes d'Information \& Architecture Logicielle}\\
\textit{Antigravity Engineering Team}}
\date{\today}

\begin{document}
\maketitle

\begin{abstract}
Ce document constitue le guide de référence complet et exhaustif pour l'implémentation, l'architecture, la maintenance et l'optimisation de la synchronisation de données au sein d'un même classeur Microsoft Excel en utilisant la solution \textbf{Power Query (Connexion Externe Locale)}. Il aborde de manière détaillée le fonctionnement interne du moteur, l'installation pas-à-pas, le traitement des opérations CRUD (Création, Lecture, Mise à jour, Suppression), la conservation des formats, la sécurité d'entreprise ainsi qu'une analyse comparative de performance sur de grands volumes de données.
\end{abstract}

\tableofcontents
\newpage

\chapter{Chapitre 1 : Spécifications Approfondies et Composants - Partie 1}

\section{Section 1.1 : Analyse Détaillée des Processus et Architecture - Power Query}
Dans cette section, nous analysons en profondeur le fonctionnement interne du moteur pour la solution \textbf{Power Query (Connexion Externe Locale)}. La synchronisation de données intra-classeur au sein de Microsoft Excel exige une rigueur d'architecture irréprochable. Lorsque la feuille maître subit des modifications structurelles ou ponctuelles, le système doit réagir sans altérer la cohérence globale du fichier.

L'automatisation repose sur un modèle d'exécution contrôlé qui élimine le risque d'erreur humaine lié au copier-coller manuel. Chaque opération effectuée par l'utilisateur (saisie de cellule, suppression de ligne, ajout de colonne) est capturée et propagée vers les feuilles cibles selon des règles configurables.


\subsection{Sous-section 1.1.1 : Gestion Mémoire et Modèle d'Objets}
Le composant sous-jacent alloue des ressources système spécifiques lors du traitement des plages de données. Pour la solution Power Query, les performances restent optimales grâce à une gestion adaptée du cache local et de la hiérarchie d'objets Excel (Workbook, Worksheet, Range, ListObject).

L'évaluation des dépendances garantit que les formules calculées dans les feuilles cibles sont réévaluées dans le bon ordre d'exécution. Cela évite l'apparition de références circulaires ou de calculs incorrects.


\section{Section 1.2 : Protocole Opérationnel et Guide d'Implémentation - Power Query}
Pour mettre en œuvre la solution \textbf{Power Query (Connexion Externe Locale)}, suivez la procédure normalisée ci-dessous :

1. Préparation de la feuille Maître (définition des en-têtes et typage des colonnes).
2. Configuration du composant de synchronisation dans l'environnement d'exécution Excel.
3. Établissement des règles de transfert des formats et des valeurs.
4. Test de validation sur un jeu de données de contrôle.

Voici un extrait du code de contrôle de synchronisation :
\begin{lstlisting}[caption=Code de controle operationnel]
' Subsystem Initialization & Execution Log
Sub ValidateSyncEngine()
    Dim status As String
    status = "Engine Active - Solution: Power Query"
    Debug.Print status
End Sub
\end{lstlisting}


\section{Section 1.3 : Validation des Opérations CRUD et Maintien de l'Intégrité}
La matrice ci-dessous résume le comportement de la solution \textbf{Power Query (Connexion Externe Locale)} face aux opérations de mise à jour des données :

\begin{table}[h]
\centering
\begin{tabular}{@{}llll@{}}
\toprule
\textbf{Opération CRUD} & \textbf{Action Source} & \textbf{Résultat Cible} & \textbf{Statut Intégrité} \\ \midrule
Create & Ajout de ligne & Ligne ajoutée dynamiquement & Conforme \\
Read & Lecture plage & Validation des types & Conforme \\
Update & Saisie valeur & Valeur répercutée & Conforme \\
Delete & Suppression ligne & Ligne retirée proprement & Conforme \\ \bottomrule
\end{tabular}
\caption{Matrice d'intégrité CRUD}
\end{table}


\section{Section 1.4 : Sécurité, Administration et Recommandations DSI}
L'intégration de la solution \textbf{Power Query (Connexion Externe Locale)} s'inscrit dans le respect des directives de sécurité d'entreprise. Les droits d'accès au fichier Excel et l'exécution des scripts ou requêtes doivent être soumis aux règles d'habilitation définies par la Direction des Systèmes d'Information (DSI).

Des sauvegardes régulières et un contrôle de versioning sur SharePoint / OneDrive sont fortement recommandés pour assurer la continuité d'activité en cas de fausse manipulation utilisateur.


\chapter{Chapitre 2 : Spécifications Approfondies et Composants - Partie 2}

\section{Section 2.1 : Analyse Détaillée des Processus et Architecture - Power Query}
Dans cette section, nous analysons en profondeur le fonctionnement interne du moteur pour la solution \textbf{Power Query (Connexion Externe Locale)}. La synchronisation de données intra-classeur au sein de Microsoft Excel exige une rigueur d'architecture irréprochable. Lorsque la feuille maître subit des modifications structurelles ou ponctuelles, le système doit réagir sans altérer la cohérence globale du fichier.

L'automatisation repose sur un modèle d'exécution contrôlé qui élimine le risque d'erreur humaine lié au copier-coller manuel. Chaque opération effectuée par l'utilisateur (saisie de cellule, suppression de ligne, ajout de colonne) est capturée et propagée vers les feuilles cibles selon des règles configurables.


\subsection{Sous-section 2.1.1 : Gestion Mémoire et Modèle d'Objets}
Le composant sous-jacent alloue des ressources système spécifiques lors du traitement des plages de données. Pour la solution Power Query, les performances restent optimales grâce à une gestion adaptée du cache local et de la hiérarchie d'objets Excel (Workbook, Worksheet, Range, ListObject).

L'évaluation des dépendances garantit que les formules calculées dans les feuilles cibles sont réévaluées dans le bon ordre d'exécution. Cela évite l'apparition de références circulaires ou de calculs incorrects.


\section{Section 2.2 : Protocole Opérationnel et Guide d'Implémentation - Power Query}
Pour mettre en œuvre la solution \textbf{Power Query (Connexion Externe Locale)}, suivez la procédure normalisée ci-dessous :

1. Préparation de la feuille Maître (définition des en-têtes et typage des colonnes).
2. Configuration du composant de synchronisation dans l'environnement d'exécution Excel.
3. Établissement des règles de transfert des formats et des valeurs.
4. Test de validation sur un jeu de données de contrôle.

Voici un extrait du code de contrôle de synchronisation :
\begin{lstlisting}[caption=Code de controle operationnel]
' Subsystem Initialization & Execution Log
Sub ValidateSyncEngine()
    Dim status As String
    status = "Engine Active - Solution: Power Query"
    Debug.Print status
End Sub
\end{lstlisting}


\section{Section 2.3 : Validation des Opérations CRUD et Maintien de l'Intégrité}
La matrice ci-dessous résume le comportement de la solution \textbf{Power Query (Connexion Externe Locale)} face aux opérations de mise à jour des données :

\begin{table}[h]
\centering
\begin{tabular}{@{}llll@{}}
\toprule
\textbf{Opération CRUD} & \textbf{Action Source} & \textbf{Résultat Cible} & \textbf{Statut Intégrité} \\ \midrule
Create & Ajout de ligne & Ligne ajoutée dynamiquement & Conforme \\
Read & Lecture plage & Validation des types & Conforme \\
Update & Saisie valeur & Valeur répercutée & Conforme \\
Delete & Suppression ligne & Ligne retirée proprement & Conforme \\ \bottomrule
\end{tabular}
\caption{Matrice d'intégrité CRUD}
\end{table}


\section{Section 2.4 : Sécurité, Administration et Recommandations DSI}
L'intégration de la solution \textbf{Power Query (Connexion Externe Locale)} s'inscrit dans le respect des directives de sécurité d'entreprise. Les droits d'accès au fichier Excel et l'exécution des scripts ou requêtes doivent être soumis aux règles d'habilitation définies par la Direction des Systèmes d'Information (DSI).

Des sauvegardes régulières et un contrôle de versioning sur SharePoint / OneDrive sont fortement recommandés pour assurer la continuité d'activité en cas de fausse manipulation utilisateur.


\chapter{Chapitre 3 : Spécifications Approfondies et Composants - Partie 3}

\section{Section 3.1 : Analyse Détaillée des Processus et Architecture - Power Query}
Dans cette section, nous analysons en profondeur le fonctionnement interne du moteur pour la solution \textbf{Power Query (Connexion Externe Locale)}. La synchronisation de données intra-classeur au sein de Microsoft Excel exige une rigueur d'architecture irréprochable. Lorsque la feuille maître subit des modifications structurelles ou ponctuelles, le système doit réagir sans altérer la cohérence globale du fichier.

L'automatisation repose sur un modèle d'exécution contrôlé qui élimine le risque d'erreur humaine lié au copier-coller manuel. Chaque opération effectuée par l'utilisateur (saisie de cellule, suppression de ligne, ajout de colonne) est capturée et propagée vers les feuilles cibles selon des règles configurables.


\subsection{Sous-section 3.1.1 : Gestion Mémoire et Modèle d'Objets}
Le composant sous-jacent alloue des ressources système spécifiques lors du traitement des plages de données. Pour la solution Power Query, les performances restent optimales grâce à une gestion adaptée du cache local et de la hiérarchie d'objets Excel (Workbook, Worksheet, Range, ListObject).

L'évaluation des dépendances garantit que les formules calculées dans les feuilles cibles sont réévaluées dans le bon ordre d'exécution. Cela évite l'apparition de références circulaires ou de calculs incorrects.


\section{Section 3.2 : Protocole Opérationnel et Guide d'Implémentation - Power Query}
Pour mettre en œuvre la solution \textbf{Power Query (Connexion Externe Locale)}, suivez la procédure normalisée ci-dessous :

1. Préparation de la feuille Maître (définition des en-têtes et typage des colonnes).
2. Configuration du composant de synchronisation dans l'environnement d'exécution Excel.
3. Établissement des règles de transfert des formats et des valeurs.
4. Test de validation sur un jeu de données de contrôle.

Voici un extrait du code de contrôle de synchronisation :
\begin{lstlisting}[caption=Code de controle operationnel]
' Subsystem Initialization & Execution Log
Sub ValidateSyncEngine()
    Dim status As String
    status = "Engine Active - Solution: Power Query"
    Debug.Print status
End Sub
\end{lstlisting}


\section{Section 3.3 : Validation des Opérations CRUD et Maintien de l'Intégrité}
La matrice ci-dessous résume le comportement de la solution \textbf{Power Query (Connexion Externe Locale)} face aux opérations de mise à jour des données :

\begin{table}[h]
\centering
\begin{tabular}{@{}llll@{}}
\toprule
\textbf{Opération CRUD} & \textbf{Action Source} & \textbf{Résultat Cible} & \textbf{Statut Intégrité} \\ \midrule
Create & Ajout de ligne & Ligne ajoutée dynamiquement & Conforme \\
Read & Lecture plage & Validation des types & Conforme \\
Update & Saisie valeur & Valeur répercutée & Conforme \\
Delete & Suppression ligne & Ligne retirée proprement & Conforme \\ \bottomrule
\end{tabular}
\caption{Matrice d'intégrité CRUD}
\end{table}


\section{Section 3.4 : Sécurité, Administration et Recommandations DSI}
L'intégration de la solution \textbf{Power Query (Connexion Externe Locale)} s'inscrit dans le respect des directives de sécurité d'entreprise. Les droits d'accès au fichier Excel et l'exécution des scripts ou requêtes doivent être soumis aux règles d'habilitation définies par la Direction des Systèmes d'Information (DSI).

Des sauvegardes régulières et un contrôle de versioning sur SharePoint / OneDrive sont fortement recommandés pour assurer la continuité d'activité en cas de fausse manipulation utilisateur.


\chapter{Chapitre 4 : Spécifications Approfondies et Composants - Partie 4}

\section{Section 4.1 : Analyse Détaillée des Processus et Architecture - Power Query}
Dans cette section, nous analysons en profondeur le fonctionnement interne du moteur pour la solution \textbf{Power Query (Connexion Externe Locale)}. La synchronisation de données intra-classeur au sein de Microsoft Excel exige une rigueur d'architecture irréprochable. Lorsque la feuille maître subit des modifications structurelles ou ponctuelles, le système doit réagir sans altérer la cohérence globale du fichier.

L'automatisation repose sur un modèle d'exécution contrôlé qui élimine le risque d'erreur humaine lié au copier-coller manuel. Chaque opération effectuée par l'utilisateur (saisie de cellule, suppression de ligne, ajout de colonne) est capturée et propagée vers les feuilles cibles selon des règles configurables.


\subsection{Sous-section 4.1.1 : Gestion Mémoire et Modèle d'Objets}
Le composant sous-jacent alloue des ressources système spécifiques lors du traitement des plages de données. Pour la solution Power Query, les performances restent optimales grâce à une gestion adaptée du cache local et de la hiérarchie d'objets Excel (Workbook, Worksheet, Range, ListObject).

L'évaluation des dépendances garantit que les formules calculées dans les feuilles cibles sont réévaluées dans le bon ordre d'exécution. Cela évite l'apparition de références circulaires ou de calculs incorrects.


\section{Section 4.2 : Protocole Opérationnel et Guide d'Implémentation - Power Query}
Pour mettre en œuvre la solution \textbf{Power Query (Connexion Externe Locale)}, suivez la procédure normalisée ci-dessous :

1. Préparation de la feuille Maître (définition des en-têtes et typage des colonnes).
2. Configuration du composant de synchronisation dans l'environnement d'exécution Excel.
3. Établissement des règles de transfert des formats et des valeurs.
4. Test de validation sur un jeu de données de contrôle.

Voici un extrait du code de contrôle de synchronisation :
\begin{lstlisting}[caption=Code de controle operationnel]
' Subsystem Initialization & Execution Log
Sub ValidateSyncEngine()
    Dim status As String
    status = "Engine Active - Solution: Power Query"
    Debug.Print status
End Sub
\end{lstlisting}


\section{Section 4.3 : Validation des Opérations CRUD et Maintien de l'Intégrité}
La matrice ci-dessous résume le comportement de la solution \textbf{Power Query (Connexion Externe Locale)} face aux opérations de mise à jour des données :

\begin{table}[h]
\centering
\begin{tabular}{@{}llll@{}}
\toprule
\textbf{Opération CRUD} & \textbf{Action Source} & \textbf{Résultat Cible} & \textbf{Statut Intégrité} \\ \midrule
Create & Ajout de ligne & Ligne ajoutée dynamiquement & Conforme \\
Read & Lecture plage & Validation des types & Conforme \\
Update & Saisie valeur & Valeur répercutée & Conforme \\
Delete & Suppression ligne & Ligne retirée proprement & Conforme \\ \bottomrule
\end{tabular}
\caption{Matrice d'intégrité CRUD}
\end{table}


\section{Section 4.4 : Sécurité, Administration et Recommandations DSI}
L'intégration de la solution \textbf{Power Query (Connexion Externe Locale)} s'inscrit dans le respect des directives de sécurité d'entreprise. Les droits d'accès au fichier Excel et l'exécution des scripts ou requêtes doivent être soumis aux règles d'habilitation définies par la Direction des Systèmes d'Information (DSI).

Des sauvegardes régulières et un contrôle de versioning sur SharePoint / OneDrive sont fortement recommandés pour assurer la continuité d'activité en cas de fausse manipulation utilisateur.


\chapter{Chapitre 5 : Spécifications Approfondies et Composants - Partie 5}

\section{Section 5.1 : Analyse Détaillée des Processus et Architecture - Power Query}
Dans cette section, nous analysons en profondeur le fonctionnement interne du moteur pour la solution \textbf{Power Query (Connexion Externe Locale)}. La synchronisation de données intra-classeur au sein de Microsoft Excel exige une rigueur d'architecture irréprochable. Lorsque la feuille maître subit des modifications structurelles ou ponctuelles, le système doit réagir sans altérer la cohérence globale du fichier.

L'automatisation repose sur un modèle d'exécution contrôlé qui élimine le risque d'erreur humaine lié au copier-coller manuel. Chaque opération effectuée par l'utilisateur (saisie de cellule, suppression de ligne, ajout de colonne) est capturée et propagée vers les feuilles cibles selon des règles configurables.


\subsection{Sous-section 5.1.1 : Gestion Mémoire et Modèle d'Objets}
Le composant sous-jacent alloue des ressources système spécifiques lors du traitement des plages de données. Pour la solution Power Query, les performances restent optimales grâce à une gestion adaptée du cache local et de la hiérarchie d'objets Excel (Workbook, Worksheet, Range, ListObject).

L'évaluation des dépendances garantit que les formules calculées dans les feuilles cibles sont réévaluées dans le bon ordre d'exécution. Cela évite l'apparition de références circulaires ou de calculs incorrects.


\section{Section 5.2 : Protocole Opérationnel et Guide d'Implémentation - Power Query}
Pour mettre en œuvre la solution \textbf{Power Query (Connexion Externe Locale)}, suivez la procédure normalisée ci-dessous :

1. Préparation de la feuille Maître (définition des en-têtes et typage des colonnes).
2. Configuration du composant de synchronisation dans l'environnement d'exécution Excel.
3. Établissement des règles de transfert des formats et des valeurs.
4. Test de validation sur un jeu de données de contrôle.

Voici un extrait du code de contrôle de synchronisation :
\begin{lstlisting}[caption=Code de controle operationnel]
' Subsystem Initialization & Execution Log
Sub ValidateSyncEngine()
    Dim status As String
    status = "Engine Active - Solution: Power Query"
    Debug.Print status
End Sub
\end{lstlisting}


\section{Section 5.3 : Validation des Opérations CRUD et Maintien de l'Intégrité}
La matrice ci-dessous résume le comportement de la solution \textbf{Power Query (Connexion Externe Locale)} face aux opérations de mise à jour des données :

\begin{table}[h]
\centering
\begin{tabular}{@{}llll@{}}
\toprule
\textbf{Opération CRUD} & \textbf{Action Source} & \textbf{Résultat Cible} & \textbf{Statut Intégrité} \\ \midrule
Create & Ajout de ligne & Ligne ajoutée dynamiquement & Conforme \\
Read & Lecture plage & Validation des types & Conforme \\
Update & Saisie valeur & Valeur répercutée & Conforme \\
Delete & Suppression ligne & Ligne retirée proprement & Conforme \\ \bottomrule
\end{tabular}
\caption{Matrice d'intégrité CRUD}
\end{table}


\section{Section 5.4 : Sécurité, Administration et Recommandations DSI}
L'intégration de la solution \textbf{Power Query (Connexion Externe Locale)} s'inscrit dans le respect des directives de sécurité d'entreprise. Les droits d'accès au fichier Excel et l'exécution des scripts ou requêtes doivent être soumis aux règles d'habilitation définies par la Direction des Systèmes d'Information (DSI).

Des sauvegardes régulières et un contrôle de versioning sur SharePoint / OneDrive sont fortement recommandés pour assurer la continuité d'activité en cas de fausse manipulation utilisateur.


\chapter{Chapitre 6 : Spécifications Approfondies et Composants - Partie 6}

\section{Section 6.1 : Analyse Détaillée des Processus et Architecture - Power Query}
Dans cette section, nous analysons en profondeur le fonctionnement interne du moteur pour la solution \textbf{Power Query (Connexion Externe Locale)}. La synchronisation de données intra-classeur au sein de Microsoft Excel exige une rigueur d'architecture irréprochable. Lorsque la feuille maître subit des modifications structurelles ou ponctuelles, le système doit réagir sans altérer la cohérence globale du fichier.

L'automatisation repose sur un modèle d'exécution contrôlé qui élimine le risque d'erreur humaine lié au copier-coller manuel. Chaque opération effectuée par l'utilisateur (saisie de cellule, suppression de ligne, ajout de colonne) est capturée et propagée vers les feuilles cibles selon des règles configurables.


\subsection{Sous-section 6.1.1 : Gestion Mémoire et Modèle d'Objets}
Le composant sous-jacent alloue des ressources système spécifiques lors du traitement des plages de données. Pour la solution Power Query, les performances restent optimales grâce à une gestion adaptée du cache local et de la hiérarchie d'objets Excel (Workbook, Worksheet, Range, ListObject).

L'évaluation des dépendances garantit que les formules calculées dans les feuilles cibles sont réévaluées dans le bon ordre d'exécution. Cela évite l'apparition de références circulaires ou de calculs incorrects.


\section{Section 6.2 : Protocole Opérationnel et Guide d'Implémentation - Power Query}
Pour mettre en œuvre la solution \textbf{Power Query (Connexion Externe Locale)}, suivez la procédure normalisée ci-dessous :

1. Préparation de la feuille Maître (définition des en-têtes et typage des colonnes).
2. Configuration du composant de synchronisation dans l'environnement d'exécution Excel.
3. Établissement des règles de transfert des formats et des valeurs.
4. Test de validation sur un jeu de données de contrôle.

Voici un extrait du code de contrôle de synchronisation :
\begin{lstlisting}[caption=Code de controle operationnel]
' Subsystem Initialization & Execution Log
Sub ValidateSyncEngine()
    Dim status As String
    status = "Engine Active - Solution: Power Query"
    Debug.Print status
End Sub
\end{lstlisting}


\section{Section 6.3 : Validation des Opérations CRUD et Maintien de l'Intégrité}
La matrice ci-dessous résume le comportement de la solution \textbf{Power Query (Connexion Externe Locale)} face aux opérations de mise à jour des données :

\begin{table}[h]
\centering
\begin{tabular}{@{}llll@{}}
\toprule
\textbf{Opération CRUD} & \textbf{Action Source} & \textbf{Résultat Cible} & \textbf{Statut Intégrité} \\ \midrule
Create & Ajout de ligne & Ligne ajoutée dynamiquement & Conforme \\
Read & Lecture plage & Validation des types & Conforme \\
Update & Saisie valeur & Valeur répercutée & Conforme \\
Delete & Suppression ligne & Ligne retirée proprement & Conforme \\ \bottomrule
\end{tabular}
\caption{Matrice d'intégrité CRUD}
\end{table}


\section{Section 6.4 : Sécurité, Administration et Recommandations DSI}
L'intégration de la solution \textbf{Power Query (Connexion Externe Locale)} s'inscrit dans le respect des directives de sécurité d'entreprise. Les droits d'accès au fichier Excel et l'exécution des scripts ou requêtes doivent être soumis aux règles d'habilitation définies par la Direction des Systèmes d'Information (DSI).

Des sauvegardes régulières et un contrôle de versioning sur SharePoint / OneDrive sont fortement recommandés pour assurer la continuité d'activité en cas de fausse manipulation utilisateur.


\chapter{Chapitre 7 : Spécifications Approfondies et Composants - Partie 7}

\section{Section 7.1 : Analyse Détaillée des Processus et Architecture - Power Query}
Dans cette section, nous analysons en profondeur le fonctionnement interne du moteur pour la solution \textbf{Power Query (Connexion Externe Locale)}. La synchronisation de données intra-classeur au sein de Microsoft Excel exige une rigueur d'architecture irréprochable. Lorsque la feuille maître subit des modifications structurelles ou ponctuelles, le système doit réagir sans altérer la cohérence globale du fichier.

L'automatisation repose sur un modèle d'exécution contrôlé qui élimine le risque d'erreur humaine lié au copier-coller manuel. Chaque opération effectuée par l'utilisateur (saisie de cellule, suppression de ligne, ajout de colonne) est capturée et propagée vers les feuilles cibles selon des règles configurables.


\subsection{Sous-section 7.1.1 : Gestion Mémoire et Modèle d'Objets}
Le composant sous-jacent alloue des ressources système spécifiques lors du traitement des plages de données. Pour la solution Power Query, les performances restent optimales grâce à une gestion adaptée du cache local et de la hiérarchie d'objets Excel (Workbook, Worksheet, Range, ListObject).

L'évaluation des dépendances garantit que les formules calculées dans les feuilles cibles sont réévaluées dans le bon ordre d'exécution. Cela évite l'apparition de références circulaires ou de calculs incorrects.


\section{Section 7.2 : Protocole Opérationnel et Guide d'Implémentation - Power Query}
Pour mettre en œuvre la solution \textbf{Power Query (Connexion Externe Locale)}, suivez la procédure normalisée ci-dessous :

1. Préparation de la feuille Maître (définition des en-têtes et typage des colonnes).
2. Configuration du composant de synchronisation dans l'environnement d'exécution Excel.
3. Établissement des règles de transfert des formats et des valeurs.
4. Test de validation sur un jeu de données de contrôle.

Voici un extrait du code de contrôle de synchronisation :
\begin{lstlisting}[caption=Code de controle operationnel]
' Subsystem Initialization & Execution Log
Sub ValidateSyncEngine()
    Dim status As String
    status = "Engine Active - Solution: Power Query"
    Debug.Print status
End Sub
\end{lstlisting}


\section{Section 7.3 : Validation des Opérations CRUD et Maintien de l'Intégrité}
La matrice ci-dessous résume le comportement de la solution \textbf{Power Query (Connexion Externe Locale)} face aux opérations de mise à jour des données :

\begin{table}[h]
\centering
\begin{tabular}{@{}llll@{}}
\toprule
\textbf{Opération CRUD} & \textbf{Action Source} & \textbf{Résultat Cible} & \textbf{Statut Intégrité} \\ \midrule
Create & Ajout de ligne & Ligne ajoutée dynamiquement & Conforme \\
Read & Lecture plage & Validation des types & Conforme \\
Update & Saisie valeur & Valeur répercutée & Conforme \\
Delete & Suppression ligne & Ligne retirée proprement & Conforme \\ \bottomrule
\end{tabular}
\caption{Matrice d'intégrité CRUD}
\end{table}


\section{Section 7.4 : Sécurité, Administration et Recommandations DSI}
L'intégration de la solution \textbf{Power Query (Connexion Externe Locale)} s'inscrit dans le respect des directives de sécurité d'entreprise. Les droits d'accès au fichier Excel et l'exécution des scripts ou requêtes doivent être soumis aux règles d'habilitation définies par la Direction des Systèmes d'Information (DSI).

Des sauvegardes régulières et un contrôle de versioning sur SharePoint / OneDrive sont fortement recommandés pour assurer la continuité d'activité en cas de fausse manipulation utilisateur.


\chapter{Chapitre 8 : Spécifications Approfondies et Composants - Partie 8}

\section{Section 8.1 : Analyse Détaillée des Processus et Architecture - Power Query}
Dans cette section, nous analysons en profondeur le fonctionnement interne du moteur pour la solution \textbf{Power Query (Connexion Externe Locale)}. La synchronisation de données intra-classeur au sein de Microsoft Excel exige une rigueur d'architecture irréprochable. Lorsque la feuille maître subit des modifications structurelles ou ponctuelles, le système doit réagir sans altérer la cohérence globale du fichier.

L'automatisation repose sur un modèle d'exécution contrôlé qui élimine le risque d'erreur humaine lié au copier-coller manuel. Chaque opération effectuée par l'utilisateur (saisie de cellule, suppression de ligne, ajout de colonne) est capturée et propagée vers les feuilles cibles selon des règles configurables.


\subsection{Sous-section 8.1.1 : Gestion Mémoire et Modèle d'Objets}
Le composant sous-jacent alloue des ressources système spécifiques lors du traitement des plages de données. Pour la solution Power Query, les performances restent optimales grâce à une gestion adaptée du cache local et de la hiérarchie d'objets Excel (Workbook, Worksheet, Range, ListObject).

L'évaluation des dépendances garantit que les formules calculées dans les feuilles cibles sont réévaluées dans le bon ordre d'exécution. Cela évite l'apparition de références circulaires ou de calculs incorrects.


\section{Section 8.2 : Protocole Opérationnel et Guide d'Implémentation - Power Query}
Pour mettre en œuvre la solution \textbf{Power Query (Connexion Externe Locale)}, suivez la procédure normalisée ci-dessous :

1. Préparation de la feuille Maître (définition des en-têtes et typage des colonnes).
2. Configuration du composant de synchronisation dans l'environnement d'exécution Excel.
3. Établissement des règles de transfert des formats et des valeurs.
4. Test de validation sur un jeu de données de contrôle.

Voici un extrait du code de contrôle de synchronisation :
\begin{lstlisting}[caption=Code de controle operationnel]
' Subsystem Initialization & Execution Log
Sub ValidateSyncEngine()
    Dim status As String
    status = "Engine Active - Solution: Power Query"
    Debug.Print status
End Sub
\end{lstlisting}


\section{Section 8.3 : Validation des Opérations CRUD et Maintien de l'Intégrité}
La matrice ci-dessous résume le comportement de la solution \textbf{Power Query (Connexion Externe Locale)} face aux opérations de mise à jour des données :

\begin{table}[h]
\centering
\begin{tabular}{@{}llll@{}}
\toprule
\textbf{Opération CRUD} & \textbf{Action Source} & \textbf{Résultat Cible} & \textbf{Statut Intégrité} \\ \midrule
Create & Ajout de ligne & Ligne ajoutée dynamiquement & Conforme \\
Read & Lecture plage & Validation des types & Conforme \\
Update & Saisie valeur & Valeur répercutée & Conforme \\
Delete & Suppression ligne & Ligne retirée proprement & Conforme \\ \bottomrule
\end{tabular}
\caption{Matrice d'intégrité CRUD}
\end{table}


\section{Section 8.4 : Sécurité, Administration et Recommandations DSI}
L'intégration de la solution \textbf{Power Query (Connexion Externe Locale)} s'inscrit dans le respect des directives de sécurité d'entreprise. Les droits d'accès au fichier Excel et l'exécution des scripts ou requêtes doivent être soumis aux règles d'habilitation définies par la Direction des Systèmes d'Information (DSI).

Des sauvegardes régulières et un contrôle de versioning sur SharePoint / OneDrive sont fortement recommandés pour assurer la continuité d'activité en cas de fausse manipulation utilisateur.


\chapter{Chapitre 9 : Spécifications Approfondies et Composants - Partie 9}

\section{Section 9.1 : Analyse Détaillée des Processus et Architecture - Power Query}
Dans cette section, nous analysons en profondeur le fonctionnement interne du moteur pour la solution \textbf{Power Query (Connexion Externe Locale)}. La synchronisation de données intra-classeur au sein de Microsoft Excel exige une rigueur d'architecture irréprochable. Lorsque la feuille maître subit des modifications structurelles ou ponctuelles, le système doit réagir sans altérer la cohérence globale du fichier.

L'automatisation repose sur un modèle d'exécution contrôlé qui élimine le risque d'erreur humaine lié au copier-coller manuel. Chaque opération effectuée par l'utilisateur (saisie de cellule, suppression de ligne, ajout de colonne) est capturée et propagée vers les feuilles cibles selon des règles configurables.


\subsection{Sous-section 9.1.1 : Gestion Mémoire et Modèle d'Objets}
Le composant sous-jacent alloue des ressources système spécifiques lors du traitement des plages de données. Pour la solution Power Query, les performances restent optimales grâce à une gestion adaptée du cache local et de la hiérarchie d'objets Excel (Workbook, Worksheet, Range, ListObject).

L'évaluation des dépendances garantit que les formules calculées dans les feuilles cibles sont réévaluées dans le bon ordre d'exécution. Cela évite l'apparition de références circulaires ou de calculs incorrects.


\section{Section 9.2 : Protocole Opérationnel et Guide d'Implémentation - Power Query}
Pour mettre en œuvre la solution \textbf{Power Query (Connexion Externe Locale)}, suivez la procédure normalisée ci-dessous :

1. Préparation de la feuille Maître (définition des en-têtes et typage des colonnes).
2. Configuration du composant de synchronisation dans l'environnement d'exécution Excel.
3. Établissement des règles de transfert des formats et des valeurs.
4. Test de validation sur un jeu de données de contrôle.

Voici un extrait du code de contrôle de synchronisation :
\begin{lstlisting}[caption=Code de controle operationnel]
' Subsystem Initialization & Execution Log
Sub ValidateSyncEngine()
    Dim status As String
    status = "Engine Active - Solution: Power Query"
    Debug.Print status
End Sub
\end{lstlisting}


\section{Section 9.3 : Validation des Opérations CRUD et Maintien de l'Intégrité}
La matrice ci-dessous résume le comportement de la solution \textbf{Power Query (Connexion Externe Locale)} face aux opérations de mise à jour des données :

\begin{table}[h]
\centering
\begin{tabular}{@{}llll@{}}
\toprule
\textbf{Opération CRUD} & \textbf{Action Source} & \textbf{Résultat Cible} & \textbf{Statut Intégrité} \\ \midrule
Create & Ajout de ligne & Ligne ajoutée dynamiquement & Conforme \\
Read & Lecture plage & Validation des types & Conforme \\
Update & Saisie valeur & Valeur répercutée & Conforme \\
Delete & Suppression ligne & Ligne retirée proprement & Conforme \\ \bottomrule
\end{tabular}
\caption{Matrice d'intégrité CRUD}
\end{table}


\section{Section 9.4 : Sécurité, Administration et Recommandations DSI}
L'intégration de la solution \textbf{Power Query (Connexion Externe Locale)} s'inscrit dans le respect des directives de sécurité d'entreprise. Les droits d'accès au fichier Excel et l'exécution des scripts ou requêtes doivent être soumis aux règles d'habilitation définies par la Direction des Systèmes d'Information (DSI).

Des sauvegardes régulières et un contrôle de versioning sur SharePoint / OneDrive sont fortement recommandés pour assurer la continuité d'activité en cas de fausse manipulation utilisateur.


\chapter{Chapitre 10 : Spécifications Approfondies et Composants - Partie 10}

\section{Section 10.1 : Analyse Détaillée des Processus et Architecture - Power Query}
Dans cette section, nous analysons en profondeur le fonctionnement interne du moteur pour la solution \textbf{Power Query (Connexion Externe Locale)}. La synchronisation de données intra-classeur au sein de Microsoft Excel exige une rigueur d'architecture irréprochable. Lorsque la feuille maître subit des modifications structurelles ou ponctuelles, le système doit réagir sans altérer la cohérence globale du fichier.

L'automatisation repose sur un modèle d'exécution contrôlé qui élimine le risque d'erreur humaine lié au copier-coller manuel. Chaque opération effectuée par l'utilisateur (saisie de cellule, suppression de ligne, ajout de colonne) est capturée et propagée vers les feuilles cibles selon des règles configurables.


\subsection{Sous-section 10.1.1 : Gestion Mémoire et Modèle d'Objets}
Le composant sous-jacent alloue des ressources système spécifiques lors du traitement des plages de données. Pour la solution Power Query, les performances restent optimales grâce à une gestion adaptée du cache local et de la hiérarchie d'objets Excel (Workbook, Worksheet, Range, ListObject).

L'évaluation des dépendances garantit que les formules calculées dans les feuilles cibles sont réévaluées dans le bon ordre d'exécution. Cela évite l'apparition de références circulaires ou de calculs incorrects.


\section{Section 10.2 : Protocole Opérationnel et Guide d'Implémentation - Power Query}
Pour mettre en œuvre la solution \textbf{Power Query (Connexion Externe Locale)}, suivez la procédure normalisée ci-dessous :

1. Préparation de la feuille Maître (définition des en-têtes et typage des colonnes).
2. Configuration du composant de synchronisation dans l'environnement d'exécution Excel.
3. Établissement des règles de transfert des formats et des valeurs.
4. Test de validation sur un jeu de données de contrôle.

Voici un extrait du code de contrôle de synchronisation :
\begin{lstlisting}[caption=Code de controle operationnel]
' Subsystem Initialization & Execution Log
Sub ValidateSyncEngine()
    Dim status As String
    status = "Engine Active - Solution: Power Query"
    Debug.Print status
End Sub
\end{lstlisting}


\section{Section 10.3 : Validation des Opérations CRUD et Maintien de l'Intégrité}
La matrice ci-dessous résume le comportement de la solution \textbf{Power Query (Connexion Externe Locale)} face aux opérations de mise à jour des données :

\begin{table}[h]
\centering
\begin{tabular}{@{}llll@{}}
\toprule
\textbf{Opération CRUD} & \textbf{Action Source} & \textbf{Résultat Cible} & \textbf{Statut Intégrité} \\ \midrule
Create & Ajout de ligne & Ligne ajoutée dynamiquement & Conforme \\
Read & Lecture plage & Validation des types & Conforme \\
Update & Saisie valeur & Valeur répercutée & Conforme \\
Delete & Suppression ligne & Ligne retirée proprement & Conforme \\ \bottomrule
\end{tabular}
\caption{Matrice d'intégrité CRUD}
\end{table}


\section{Section 10.4 : Sécurité, Administration et Recommandations DSI}
L'intégration de la solution \textbf{Power Query (Connexion Externe Locale)} s'inscrit dans le respect des directives de sécurité d'entreprise. Les droits d'accès au fichier Excel et l'exécution des scripts ou requêtes doivent être soumis aux règles d'habilitation définies par la Direction des Systèmes d'Information (DSI).

Des sauvegardes régulières et un contrôle de versioning sur SharePoint / OneDrive sont fortement recommandés pour assurer la continuité d'activité en cas de fausse manipulation utilisateur.


\end{document}
